\documentclass[11pt,a4paper]{article}

\usepackage{amsmath,amssymb}
\usepackage{physics}
\usepackage{hyperref}
\usepackage[margin=2.5cm]{geometry}
\usepackage{booktabs}

\newcommand{\code}[1]{\texttt{#1}}
\newcommand{\Trsc}{\operatorname{Tr}_{sc}}

\title{\textbf{EMT Disconnected One-Point Building Blocks}\\[4pt]
       \large Flowed Quark/Gluon Loops, HDF5 Metadata, and Hierarchical Probing}
\author{PyQUDA Measurement Notes}
\date{\today}

\begin{document}
\maketitle
\tableofcontents

\section{Purpose}

This note documents the hadron-independent one-point functions used to build
disconnected contributions to pion and proton energy-momentum tensor (EMT)
matrix elements.  The corresponding source file is
\code{pyquda\_measurement\_utils/Disconnected\_1pt\_EMT\_vibe\_develop.py}.

The one-point workflows do not form a hadron three-point function by themselves.
Instead, they provide flowed quark and gluon loop observables that are combined
with hadron two-point functions in ensemble analysis.

\section{Flowed One-Point Observables}

Let \(t_f\) denote the gradient-flow time.  The flowed gluonic building block is
schematically
\begin{equation}
  L^g_{\mu\nu}(q,t_f)
  =
  \sum_{\mathbf{x}}
  e^{i\mathbf{q}\cdot\mathbf{x}}\,
  O^g_{\mu\nu}(\mathbf{x},\tau;t_f),
\end{equation}
where \(O^g_{\mu\nu}\) is the chosen flowed gluonic EMT operator.

The flowed quark loop is estimated stochastically.  For a noise vector
\(\xi_r\) and solution \(\eta_r = D^{-1}\xi_r\), the measured loop has the
form
\begin{equation}
  L^q_{\mu\nu}(q,\tau;t_f)
  =
  \frac{1}{N_r}
  \sum_{r=1}^{N_r}
  \sum_{\mathbf{x}}
  e^{i\mathbf{q}\cdot\mathbf{x}}\,
  \Trsc\!\left[
    \xi_r^\dagger(x,t_f)\,
    \Gamma_{\mu\nu}[U(t_f),D]\,
    \eta_r(x,t_f)
  \right].
\end{equation}
The operator \(\Gamma_{\mu\nu}\) includes the symmetrized covariant derivative
structure used by the EMT kernel.  The code stores only the upper triangle
\(\mu \leq \nu\), and downstream analysis reconstructs the lower triangle by
symmetry when needed.

\section{Disconnected Combination with Hadron Two-Point Functions}

For a hadron two-point function \(C_2^H(p_f,t)\), the disconnected three-point
building block is
\begin{equation}
  C^{H,\mathrm{disc}}_{3,\mu\nu}
  (p_f,p_i;t,\tau;t_f)
  =
  \left\langle
    C_2^H(p_f,t)\,
    L_{\mu\nu}(q,\tau;t_f)
  \right\rangle_{\mathrm{cfg}}
  -
  \left\langle C_2^H(p_f,t)\right\rangle_{\mathrm{cfg}}
  \left\langle L_{\mu\nu}(q,\tau;t_f)\right\rangle_{\mathrm{cfg}},
\end{equation}
with
\begin{equation}
  q = p_f - p_i.
\end{equation}
The associated ratio is typically
\begin{equation}
  R^{H,\mathrm{disc}}_{\mu\nu}(t,\tau;t_f)
  =
  \frac{
    C^{H,\mathrm{disc}}_{3,\mu\nu}(p_f,p_i;t,\tau;t_f)
  }{
    C_2^H(p_f,t)
  },
\end{equation}
before applying kinematic projections, vacuum-subtracted gradient-flow matching,
and the perturbative mixing coefficients required for the renormalized EMT.

The subtraction term is essential: the one-point loop is hadron independent and
contains a vacuum expectation value.  The subtraction should be done at the
ensemble-analysis level using the same gauge-configuration sample as the hadron
two-point function.

\section{Ringed-Fermion Diagnostics}

The quark workflow also writes \(\chi\)-type flowed-fermion diagnostics.  These
outputs are useful for monitoring or reconstructing ringed-fermion normalization
factors.  The standard kinetic normalization is not the raw \(\chi\) dataset by
itself; it should be reconstructed from the appropriate flowed kinetic
expectation values, especially the zero-momentum diagonal EMT components.

\section{Hierarchical Probing}

The quark loop supports two source schemes:
\begin{equation}
  \code{noise\_scheme} \in
  \{\code{zn},\code{hierarchical\_probing}\}.
\end{equation}
For ordinary \(Z_n\) noise, each stochastic source corresponds to one inversion.
For hierarchical probing, a base noise vector \(\xi_b\) is multiplied by a set of
site-only Rademacher vectors \(h_k(x)=\pm 1\):
\begin{equation}
  \xi_{b,k}(x) = h_k(x)\,\xi_b(x).
\end{equation}
The effective number of inversions is therefore
\begin{equation}
  N_{\mathrm{eff}}
  =
  N_{\mathrm{base}}\,
  N_{\mathrm{HP}}.
\end{equation}
In the application this is
\begin{equation}
  \code{effective\_n\_inversions}
  =
  \code{EMT\_1PT\_N\_VEC}
  \times
  \code{EMT\_1PT\_HP\_NUM\_VECTORS}.
\end{equation}
\code{EMT\_1PT\_HP\_NUM\_VECTORS} must be a positive power of two.

\subsection{Ordering Choices}

Two Gray-code ordering choices are implemented:
\begin{center}
\begin{tabular}{ll}
\toprule
Ordering & Meaning \\
\midrule
\code{global\_xyzt\_gray\_projected\_to\_evenodd}
& Gray code the full \(x,y,z,t\) site id. \\
\code{spatial\_xyz\_then\_t\_gray\_projected\_to\_evenodd}
& Gray code spatial sites first, then append time bits. \\
\bottomrule
\end{tabular}
\end{center}

The global ordering is the default to preserve the previously validated
behavior.  The spatial-first ordering is intended for studies where the leading
noise correlations are expected to be dominated by spatial locality and where
time dilution is not yet being used.

Spin-color dilution and time dilution are not part of the current short-term
workflow.

\section{HDF5 Output}

The quark HDF5 file stores metadata attributes:
\begin{center}
\begin{tabular}{ll}
\toprule
Attribute & Meaning \\
\midrule
\code{flow\_type}, \code{flow\_epsilon}, \code{flow\_steps}, \code{flow\_times}
& Gradient-flow schedule. \\
\code{qext}, \code{pf}, \code{p\_2pt}
& Momentum labels used by downstream analysis. \\
\code{volume\_norm}
& Spatial volume normalization used in the averaged datasets. \\
\code{upper\_triangle\_only}
& Only \(\mu \leq \nu\) components are written. \\
\code{noise\_scheme}, \code{hp\_num\_vectors}, \code{hp\_ordering}
& Source-generation scheme. \\
\code{n\_base\_noise}, \code{effective\_n\_inversions}
& Stochastic source counting. \\
\bottomrule
\end{tabular}
\end{center}

The raw quark datasets are:
\begin{equation}
  \code{raw/Tmunu\_pervec},\quad
  \code{raw/CHI\_pervec},\quad
  \code{raw/source\_index},\quad
  \code{raw/base\_noise\_index},\quad
  \code{raw/hp\_index}.
\end{equation}
The bookkeeping datasets allow analysis code to group effective HP sources back
to the original base noise:
\begin{center}
\begin{tabular}{ll}
\toprule
Dataset & Meaning \\
\midrule
\code{source\_index}
& Effective source index after HP expansion. \\
\code{base\_noise\_index}
& Original stochastic base-noise index. \\
\code{hp\_index}
& Hierarchical-probing vector index for that base noise. \\
\bottomrule
\end{tabular}
\end{center}

\section{Current Validation Status}

The current implementation has been checked with:
\begin{itemize}
  \item Python syntax compilation.
  \item A CPU/numpy backend check that \(\mathrm{HP}=1\) leaves the base source
        unchanged.
  \item A Perlmutter login-node GPU smoke test with
        \(\code{hp\_num\_vectors}=2\) and spatial-first ordering.
\end{itemize}

Before production analysis, the recommended next validation is a scaling study
comparing \(Z_n\), \(\mathrm{HP}=1\), \(\mathrm{HP}=2,4,8,16\) at fixed gauge,
flow time, and momentum.

\end{document}
